% Part: first-order-logic
% Chapter: models-theories
% Section: introduction

\documentclass[../../../include/open-logic-section]{subfiles}

\begin{document}

\olfileid{fol}{mat}{int}
\olsection{مقدمه}

\begin{explain}
پرورش روش اصل‌موضوعی دستاوردی مهم در تاریخ علم و، به‌ویژه، در تاریخ
ریاضیات است. پرورش اصل‌موضوعیِ یک حوزه مستلزم روشن‌ساختن پرسش‌های
بسیاری است: این حوزه دربارهٔ چیست؟ بنیادی‌ترین مفاهیم آن کدام‌اند؟
این مفاهیم چه نسبتی با یکدیگر دارند؟ آیا می‌توان همهٔ مفاهیم حوزه را
بر حسب این مفاهیم بنیادی تعریف کرد؟ این مفاهیم از چه قوانینی پیروی
می‌کنند و باید پیروی کنند؟

روش اصل‌موضوعی و منطق گویی برای یکدیگر ساخته شده‌اند. منطق صوری ابزار
صورت‌بندی نظریه‌های اصل‌موضوعی، اثبات قضیه‌ها از اصول موضوع نظریه به
شیوه‌ای دقیقاً معین، و بررسی نظام‌مند خواص همهٔ دستگاه‌هایی را فراهم
می‌کند که اصول موضوع را برآورده می‌کنند.
\end{explain}

\begin{defn}
مجموعه‌ای چون~$\Gamma$ از !!{sentence}s \emph{بسته} است هرگاه و تنها
هرگاه، اگر $\Gamma \Entails !A$ آنگاه $!A \in \Gamma$. \emph{بستار}
مجموعه‌ای چون~$\Gamma$ از !!{sentence}s عبارت است از
$\Setabs{!A}{\Gamma \Entails !A}$.

می‌گوییم~$\Gamma$ به‌وسیلهٔ مجموعه‌ای چون~$\Delta$ از جمله‌ها
\emph{اصل‌موضوع‌بندی شده است} اگر $\Gamma$ بستار~$\Delta$ باشد.
\end{defn}

\begin{explain}
می‌توانیم نظریه‌ای اصل‌موضوعی را مجموعهٔ جمله‌هایی بدانیم که
مجموعهٔ اصول موضوعش~$\Delta$ آن را اصل‌موضوع‌بندی می‌کند. به بیان دیگر،
فرض کنید زبانی مرتبهٔ اول داریم که شامل نمادهای غیرمنطقی برای مفاهیم
اولیهٔ علمی است که می‌خواهیم به شیوهٔ اصل‌موضوعی بررسی کنیم، و نیز
مجموعه‌ای از !!{sentence}s داریم که قوانین بنیادی آن علم را بیان
می‌کنند. در این صورت می‌توانیم نظریه را با همهٔ !!{sentence}s این زبان
که از اصول موضوع نتیجه می‌شوند بازنمایی کنیم. دامنهٔ نمونه‌ها از موارد
ساده‌ای با تنها یک مفهوم اولیه و اصول موضوع ساده، مانند نظریهٔ
ترتیب‌های جزئی، تا نظریه‌های پیچیده‌ای چون مکانیک نیوتنی گسترده است.

واقعیت‌های منطقی مهمی که این رویکرد صوری را به روش اصل‌موضوعی چنین
بااهمیت می‌سازند از قرار زیرند. فرض کنید $\Gamma$ دستگاهی از اصول موضوع
برای یک نظریه، یعنی مجموعه‌ای از جمله‌ها، باشد.
\begin{enumerate}
\item می‌توانیم به‌دقت بیان کنیم که یک دستگاه اصول موضوع چه هنگام ردهٔ
  مقصودی از !!{structure}s را مشخصه‌سازی می‌کند. یعنی اگر به ردهٔ معینی
  از !!{structure}s نظر داشته باشیم، دستگاه اصول موضوع~$\Gamma$ آن رده
  را با موفقیت مشخصه‌سازی می‌کند هرگاه و تنها هرگاه !!{structure}s آن
  رده دقیقاً همان~$\Struct M$هایی باشند که $\Sat{M}{\Gamma}$.
\item ممکن است در این کار ناکام بمانیم، زیرا $\Struct M$هایی وجود دارند
  که $\Sat{M}{\Gamma}$، اما $\Struct M$ از !!{structure}s مورد نظر ما
  نیست. این امر ممکن است ما را به افزودن اصول موضوعی وادارد که در
  در~$\Struct M$ صادق نیستند.
\item اگر دست‌کم از این حیث موفق باشیم که $\Gamma$ در همهٔ
  !!{structure}s مورد نظر صادق باشد، آنگاه هرگاه $\Gamma \Entails !A$،
  جملهٔ~$!A$ در همهٔ !!{structure}s مورد نظر صادق است. بنابراین
  می‌توانیم صرفاً با نشان‌دادن اینکه جمله‌ها از اصول موضوع نتیجه
  می‌شوند، با ابزارهای منطقی (مانند روش‌های !!{derivation}) نشان دهیم
  که آن‌ها در همهٔ !!{structure}s مورد نظر صادق‌اند.
\item گاهی هیچ !!{structure} مقصودی در نظر نداریم و در عوض از خود اصول
  موضوع آغاز می‌کنیم: کار را با چند مفهوم اولیه شروع می‌کنیم که
  می‌خواهیم قوانین معینی را برآورده کنند و این قوانین را در دستگاهی از
  اصول موضوع صورت‌بندی می‌کنیم. یکی از نخستین چیزهایی که می‌خواهیم
  وارسی کنیم این است که اصول موضوع با یکدیگر تناقض نداشته باشند: اگر
  متناقض باشند، هیچ مفهومی نمی‌تواند از این قوانین پیروی کند و کوشیده‌ایم
  نظریه‌ای نامنسجم برپا کنیم. با یافتن مدلی از~$\Gamma$ می‌توانیم
  وارسی کنیم که چنین اتفاقی نمی‌افتد. و اگر نظریهٔ ما مدل‌هایی داشته
  باشد، می‌توانیم با روش‌های منطقی آن‌ها را بررسی کنیم و نیز مدل‌هایی
  بسازیم.
\item استقلال اصول موضوع نیز پرسشی مهم است. ممکن است یکی از اصول موضوع
  در واقع پیامد اصول دیگر و ازاین‌رو زائد باشد. می‌توانیم با اثبات
  $\Gamma \setminus \{!A\} \Entails !A$ نشان دهیم که اصل موضوع~$!A$ در
  $\Gamma$ زائد است. همچنین می‌توانیم با نشان‌دادن ارضاپذیری
  $(\Gamma \setminus \{!A\}) \cup \{\lnot !A\}$ اثبات کنیم که یک اصل
  موضوع زائد نیست. برای نمونه، استقلال اصل توازی از دیگر اصول موضوع
  هندسه به همین شیوه نشان داده شد.
\item پرسش مهم دیگر تعریف‌پذیری مفاهیم در یک نظریه است: انتخاب زبان
  تعیین می‌کند مدل‌های یک نظریه از چه چیزهایی تشکیل شوند. اما لازم نیست
  هر جنبهٔ نظریه در مدل‌هایش جداگانه بازنمایی شود. برای نمونه، هر
  ترتیب~$\le$ ترتیب اکید متناظری چون~$<$ را تعیین می‌کند---اگر یکی داده
  شود، می‌توانیم دیگری را تعریف کنیم. پس لازم نیست مدل نظریه‌ای که چنین
  ترتیبی را در بر می‌گیرد، ترتیب اکید متناظر را \emph{نیز} در خود داشته
  باشد. در حالت کلی، چه هنگام می‌توان یک رابطه را بر حسب روابط دیگر
  تعریف کرد؟ چه هنگام تعریف یک رابطه بر حسب روابط دیگر ناممکن است (و
  ازاین‌رو باید آن را به مفاهیم اولیهٔ زبان افزود)؟
\end{enumerate}
\end{explain}


\end{document}
